\section{Repeatability of Command Changes across Policy-Noise Draws}
\label{sec:supp-q0-consistency}
Does a perturbation change the commands in the same way under a different draw of policy noise? For the arm it does: across the four families, the agreement of the translation changes is $0.884$--$0.920$ in $\pi_{0.5}$ and $0.799$--$0.889$ in $\pi_0$, where 1 means identical changes (Table~\ref{tab:q0-consistency-primary}). This supports Sec.~\ref*{M-sec:early} of the main paper: the first five commands are a repeatable readout of a perturbation. Gripper changes agree much less.

\paragraph{Rotation repeats nearly as well; the gripper does not.}
Rotation agreement spans $0.818$--$0.858$ in $\pi_{0.5}$ and $0.719$--$0.822$ in $\pi_0$. Agreement of the continuous gripper command is lower, $0.143$--$0.371$ and $0.077$--$0.342$. These ranges span the family point estimates; Table~\ref{tab:q0-consistency-primary} gives the intervals for each component.

\paragraph{The same perturbation under two noise draws.}
We use the commands of the paired executions (Sec.~\ref{sec:five-arm}). Within each policy--setting record, $R_0/T_0$ and $R_1/T_1$ share a policy-noise draw. All four runs must have five commands at the first query ($q=0$), with verified initial state, source, noise draw and command positions; a command stays missing if its environment call did not return. All outcome groups, including joint failure, are eligible.

\paragraph{A signed agreement score.}
The score divides the inner product of the two noise-matched changes by their mean squared size, so equal changes score one, unequal changes in the same direction score below one, and opposing changes score below zero. Flatten the first five commands and write $d_n=T_n-R_n$, $n\in\{0,1\}$; translation and rotation have 15 coordinates each and the gripper five. Define
\begin{align}
D&=\tfrac12(\|d_0\|^2+\|d_1\|^2),\\
M&=\left\|\tfrac12(d_0+d_1)\right\|^2,
\quad V=\tfrac14\|d_0-d_1\|^2,\\
U&=d_0^\top d_1=M-V,\qquad D=M+V.
\end{align}
In compact notation, with $\delta u_n=d_n$ the change in the first five commands under noise draw $n$, the numerator and denominator are
\begin{equation}
U=\delta u_0^\top\delta u_1,\quad
D=\tfrac12\bigl(\|\delta u_0\|^2+\|\delta u_1\|^2\bigr).
\label{eq:command-consistency}
\end{equation}


The agreement for one setting is
\begin{equation}
s=\frac{U}{D}
=\frac{2d_0^\top d_1}{\|d_0\|^2+\|d_1\|^2}
=\frac{2r\cos\theta}{1+r^2},\quad s\in[-1,1],
\label{eq:q0-consistency}
\end{equation}
where the last form requires both norms to be positive and $r=\|d_1\|/\|d_0\|$. If both changes are zero, $s$ is undefined; if only one is zero, $s=0$ while the cosine is undefined. Ratios use the observed changes directly, without an epsilon or clipping. Reading $U$ as a squared population mean would require independent, identically distributed noise; the ratio is a descriptive score for the two recorded draws, not an unbiased estimator of a population ratio.

\paragraph{Support and weighting.}
The primary plan has 202 settings on 28 common tasks per policy. All four runs have the first five commands in 197/202 settings for $\pi_{0.5}$ and 198/202 for $\pi_0$; Light has 43/48 settings on 26 tasks and 44/48 on 27, while the other families keep all primary settings and 28 tasks (Table~\ref{tab:q0-consistency-primary}). The full sensitivity scope keeps 258/268 and 263/268 settings (stored full-scope table); its families originally cover 34 Background, 32 Camera, 40 Light and 39 Sensor Noise tasks. Command eligibility and outcome availability are counted separately; the 39 unknown outcomes stay in the denominator of 2,680 runs.

Per-setting scores are averaged within task, and tasks equally, per policy, family and component; the estimate is a mean of ratios. The 2,000 whole-task bootstrap draws share their task resampling across series, keep settings and noise pairs nested, and are all valid; the 95\% percentile intervals are pointwise and condition on the two noise draws and the selected settings.

\begin{table*}[t]
\caption{\textbf{Translation and rotation changes repeat across policy-noise draws; the gripper agrees less.} Primary 28-task plan per policy; Settings and Tasks give defined/planned support. Entries are task-equal means of the per-setting agreement $U/D$ with 95\% whole-task intervals. These are normalized inner products of the changes under the two draws, not energy percentages.}

\centering\small
\setlength{\tabcolsep}{3.5pt}
\begin{tabular}{llccccc}
\toprule
Policy & Family & Settings & Tasks & Translation & Rotation & Continuous gripper\\
\midrule
$\pi_{0.5}$ & Background & 48/48 & 28/28 & $0.898\, [0.843,0.946]$ & $0.858\, [0.793,0.913]$ & $0.371\, [0.270,0.467]$\\
$\pi_{0.5}$ & Camera & 51/51 & 28/28 & $0.920\, [0.872,0.957]$ & $0.818\, [0.733,0.892]$ & $0.362\, [0.282,0.442]$\\
$\pi_{0.5}$ & Light & 43/48 & 26/28 & $0.884\, [0.810,0.943]$ & $0.848\, [0.778,0.907]$ & $0.143\, [0.009,0.279]$\\
$\pi_{0.5}$ & Sensor Noise & 55/55 & 28/28 & $0.909\, [0.870,0.942]$ & $0.849\, [0.788,0.908]$ & $0.371\, [0.250,0.491]$\\
\midrule
$\pi_0$ & Background & 48/48 & 28/28 & $0.865\, [0.810,0.909]$ & $0.807\, [0.720,0.877]$ & $0.342\, [0.219,0.453]$\\
$\pi_0$ & Camera & 51/51 & 28/28 & $0.889\, [0.840,0.928]$ & $0.814\, [0.756,0.866]$ & $0.325\, [0.185,0.460]$\\
$\pi_0$ & Light & 44/48 & 27/28 & $0.799\, [0.711,0.880]$ & $0.719\, [0.633,0.797]$ & $0.077\, [-0.090,0.245]$\\
$\pi_0$ & Sensor Noise & 55/55 & 28/28 & $0.844\, [0.780,0.896]$ & $0.822\, [0.762,0.876]$ & $0.228\, [0.078,0.388]$\\
\bottomrule
\end{tabular}

\label{tab:q0-consistency-primary}
\end{table*}


\paragraph{Opposing changes and zero changes stay in the comparison.}
The continuous gripper changes in opposite directions ($U<0$) in 36/197 and 62/198 primary settings, and the $\pi_0$ Light gripper interval crosses zero in both scopes. Translation never has $U<0$; $\pi_0$ Sensor Noise rotation has it once. All scores are defined and no setting has only one zero change; the rules for zero changes remain in the protocol. In the full scope the gripper has $U<0$ in 52/258 and 75/263 settings; this larger scope is a sensitivity analysis, not an independent replication.

Light has different eligible settings in the two policies, so a direct comparison of the policies requires their common settings. The gripper command can also be read as an event: Sec.~\ref{sec:supp-landmark-final} tests whether the timing of its first opening-to-closing transition adds information about lost success at a fixed cutoff.

\subsection{Gripper timing does not clearly improve prediction of lost success}
\label{sec:supp-landmark-final}
Does the timing of the first opening-to-closing gripper command add information about lost success beyond the size of the command changes? At a cutoff of 60 controls, fixed in advance, adding timing lowers the held-out Brier loss slightly in both policies, with intervals that include zero (Table~\ref{tab:landmark-final}). This is an additional, exploratory study; the main paper does not rely on it. The feature is a command event, not physical closure or object contact.

\begin{table*}[t]
\caption{\textbf{Adding gripper timing to the command-change sizes does not clearly lower the loss at a fixed cutoff of 60 controls.} Cutoff pairs (of 404 per policy) reach the cutoff with valid command records and outcomes; the subset whose clean run succeeds (N/tasks) is analyzed. Improvement is baseline minus augmented out-of-fold Brier loss, positive when timing lowers the loss, with a 97.5\% paired whole-task interval. The denominator 404 is the fixed plan per policy; 375/360 pairs are analyzed.}

\centering\small
\begin{tabular}{lrrrrr}\toprule
Policy & Cutoff pairs & Clean-success N/tasks & Baseline Brier & Augmented Brier & Improvement [97.5\% interval]\\\midrule
$\pi_{0.5}$ & 396/404 & 375/27 & 0.09664 & 0.09555 & $0.00109[-0.00224,0.00448]$\\
$\pi_0$ & 397/404 & 360/26 & 0.12323 & 0.12172 & $0.00151[-0.00471,0.00724]$\\\bottomrule
\end{tabular}

\label{tab:landmark-final}
\end{table*}

\paragraph{Event reconstruction.}
Across all 50 runs, the Panda gripper command is positive for closing and negative for opening; zero leaves the accumulated command unchanged. Events use only complete, finite command records with valid identifiers up to the cutoff; the 39 unknown run outcomes, early terminations and technical failures are kept.

\paragraph{Baseline and augmented predictors.}
Of the 404 primary pairs per policy, both runs must reach the cutoff with valid command records and outcomes, and the clean run must succeed. Both predictors use only information available at the cutoff. The baseline uses $\log(1+\mathrm{RMS})$ of the translation, rotation and continuous-gripper command differences over the same first 60 controls, plus family indicators. The augmented predictor adds whether each run has an opening-to-closing transition by the cutoff, the difference of their times divided by 60 when both do, and each run's initial command direction; occurrence flags separate an absent event from a zero time difference.

\paragraph{Fitting on held-out tasks and intervals.}
The ridge model, fixed in advance, uses penalty $\alpha=1$, an unpenalized intercept, centering and scaling on the training data only, and leave-one-task-out prediction clipped to $[0,1]$; tasks are weighted equally, and families, settings and noise draws equally within their parents. Both predictors use identical rows in each fold; no training fold has a single class (which would fall back to the training mean). Both initial-closing indicator columns are zero in all eligible rows and add no variation. Paired whole-task resampling (2,000 draws) gives 97.5\% intervals per policy; the predictors are not refitted inside the bootstrap.

\paragraph{Outcomes of the analyzed pairs.}
On the primary support the outcome counts are $328/47/17/4$ SS/SF/FS/FF with eight unknown pairs for $\pi_{0.5}$ and $309/51/10/27$ with seven unknown pairs for $\pi_0$; every validated pair is kept whether or not a transition occurs. Pairs whose clean run succeeds number $93/98/82/102$ (Background/Camera/Light/Sensor Noise) for $\pi_{0.5}$ and $88/89/81/102$ for $\pi_0$. The estimates in Table~\ref{tab:landmark-final} apply to this conditional support.

